\documentclass[11pt]{article}

\usepackage[T1]{fontenc}
\usepackage{lmodern}
\usepackage{microtype}
\usepackage[margin=1in]{geometry}
\usepackage{amsmath,amssymb,amsthm,mathtools}
\usepackage[colorlinks=true,linkcolor=blue,citecolor=blue,urlcolor=blue]{hyperref}

\newtheorem{theorem}{Theorem}
\newtheorem{lemma}[theorem]{Lemma}
\newtheorem{corollary}[theorem]{Corollary}
\newtheorem{remark}[theorem]{Remark}
\newtheorem{definition}[theorem]{Definition}

\DeclareMathOperator{\Cay}{Cay}
\newcommand{\Cdir}{\vec C}
\newcommand{\torus}[2]{D_{#1}(#2)}

\title{Composite expansion for Hamilton decompositions of directed tori}
\author{}
\date{14 March 2026}

\begin{document}
\maketitle

\begin{abstract}
We formalize a general Cartesian-power lifting principle behind the passage from
known low-dimensional directed tori to higher composite dimensions.
The key theorem is this: if a finite digraph $G$ on $n$ vertices admits a
Hamilton decomposition, and if the Cartesian power $\Cdir_n^{\square b}$ admits
one, then the Cartesian power $G^{\square b}$ admits a Hamilton decomposition as
well.
For directed tori this yields the pointwise composite-expansion implication
\[
\torus{a}{m}\text{ and }\torus{b}{m^a}
\Longrightarrow
\torus{ab}{m},
\]
and hence the multiplicative closure of the set of dimensions that are solved
for every modulus.
The previously discussed statement ``if $G$ has a Hamilton decomposition then so
does $G\square G$'' appears as the special case $b=2$, using the explicit
decomposition of $\Cdir_n\square\Cdir_n$ into two directed Hamilton cycles.
\end{abstract}

\section{Definitions}

\begin{definition}[Cartesian product and Cartesian powers]
Let $G=(V,A)$ and $H=(W,B)$ be finite digraphs.
Their \emph{Cartesian product} $G\square H$ is the digraph with vertex set
$V\times W$ and arc set consisting of all arcs
\[
(x,y)\to(x',y) \quad\text{whenever } x\to x'\text{ is an arc of }G,
\]
and all arcs
\[
(x,y)\to(x,y') \quad\text{whenever } y\to y'\text{ is an arc of }H.
\]
For $b\ge 1$, write $G^{\square b}$ for the $b$-fold Cartesian power of $G$.
\end{definition}

\begin{definition}[Hamilton decomposition]
A \emph{directed Hamilton cycle} in a finite digraph is a spanning directed
cycle.
A \emph{Hamilton decomposition} of a finite digraph is a partition of its arc
set into arc-disjoint directed Hamilton cycles.
\end{definition}

For $n\ge 2$, write $\Cdir_n$ for the directed cycle on $\mathbb Z_n$ with arcs
$t\to t+1$ modulo $n$.

\section{A Cartesian-power lift theorem}

\begin{theorem}[Cartesian-power lift]\label{thm:power-lift}
Let $G$ be a finite digraph on $n\ge 2$ vertices, and suppose that $A(G)$ is
partitioned by directed Hamilton cycles
\[
H_0,H_1,\dots,H_{k-1}.
\]
Let $b\ge 1$.
Assume that $\Cdir_n^{\square b}$ admits a Hamilton decomposition into $b$
directed Hamilton cycles.
Then $G^{\square b}$ admits a Hamilton decomposition into $kb$ directed Hamilton
cycles.
\end{theorem}

\begin{proof}
Let
\[
G^{\square b}=G^{(1)}\square G^{(2)}\square\cdots\square G^{(b)},
\]
where each $G^{(r)}$ is a copy of $G$.
For each factor $r\in\{1,\dots,b\}$ and each color index
$i\in\{0,\dots,k-1\}$, let $H_i^{(r)}$ denote the corresponding copy of $H_i$
inside $G^{(r)}$.
Define a spanning subdigraph $F_i\subseteq G^{\square b}$ by
\[
F_i:=H_i^{(1)}\square H_i^{(2)}\square\cdots\square H_i^{(b)}.
\]
Because every $H_i^{(r)}$ is a directed Hamilton cycle on $n$ vertices,
there is a digraph isomorphism
\[
F_i\cong \Cdir_n^{\square b}.
\]
By hypothesis, each $F_i$ therefore admits a Hamilton decomposition into $b$
directed Hamilton cycles.

It remains to show that the subdigraphs $F_0,\dots,F_{k-1}$ partition the arc
set of $G^{\square b}$.
An arc of $G^{\square b}$ changes exactly one coordinate.
So every arc has the form
\[
(x_1,\dots,x_b)\to(x_1,\dots,x_{r-1},x_r',x_{r+1},\dots,x_b)
\]
for a unique index $r$, where $x_r\to x_r'$ is an arc of the copy $G^{(r)}$.
Since the Hamilton cycles $H_0^{(r)},\dots,H_{k-1}^{(r)}$ partition the arc set
of $G^{(r)}$, there is a unique index $i$ such that
\[
x_r\to x_r'\in A\bigl(H_i^{(r)}\bigr).
\]
By definition of the Cartesian product, that full product arc then belongs to
$F_i$, and it belongs to no $F_j$ with $j\neq i$.
Thus
\[
A\bigl(G^{\square b}\bigr)=\bigsqcup_{i=0}^{k-1} A(F_i).
\]
Now decompose each $F_i$ into its $b$ Hamilton cycles, and collect all of those
decompositions together.
Because the $F_i$ are pairwise arc-disjoint and cover all arcs, this produces a
Hamilton decomposition of $G^{\square b}$ into $kb$ directed Hamilton cycles.
\end{proof}

\begin{remark}
The theorem is best viewed as a two-level construction.
First one colors every factor-copy of $G$ by the Hamilton decomposition of $G$.
Then one freezes a color $i$ across all $b$ factors, obtaining the monochromatic
product block $F_i\cong \Cdir_n^{\square b}$.
A Hamilton decomposition of $\Cdir_n^{\square b}$ can then be imported into each
block $F_i$ independently.
\end{remark}

\section{The square case}

The self-square statement is the special case $b=2$.
For completeness we include the explicit $2$-cycle decomposition of
$\Cdir_n\square\Cdir_n$.

\begin{lemma}[Square lemma]\label{lem:square}
For every integer $n\ge 2$, the digraph $\Cdir_n\square\Cdir_n$ decomposes into
two directed Hamilton cycles.
\end{lemma}

\begin{proof}
Identify the vertex set with $\mathbb Z_n^2$.
At each vertex $(u,v)$ there are exactly two outgoing arcs,
\[
(u,v)\to(u+1,v),
\qquad
(u,v)\to(u,v+1),
\]
where all arithmetic is modulo $n$.
Define maps $\alpha,\beta:\mathbb Z_n^2\to\mathbb Z_n^2$ by
\[
\alpha(u,v)=
\begin{cases}
(u+1,v), & u+v\not\equiv n-1 \pmod n,\\[2mm]
(u,v+1), & u+v\equiv n-1 \pmod n,
\end{cases}
\]
\[
\beta(u,v)=
\begin{cases}
(u,v+1), & u+v\not\equiv n-1 \pmod n,\\[2mm]
(u+1,v), & u+v\equiv n-1 \pmod n.
\end{cases}
\]
At each vertex, $\alpha$ and $\beta$ choose complementary outgoing arcs, so
their arc sets partition $A(\Cdir_n\square\Cdir_n)$.

For $\alpha$, define the bijection
\[
\Psi(u,v)=(u+v,\,v).
\]
Since $\Psi^{-1}(s,t)=(s-t,\,t)$, a direct substitution gives
\[
\Psi\alpha\Psi^{-1}(s,t)=
\begin{cases}
(s+1,t), & s\not\equiv n-1 \pmod n,\\[2mm]
(0,t+1), & s\equiv n-1 \pmod n.
\end{cases}
\]
Thus $\alpha$ is conjugate to the standard odometer on $\mathbb Z_n^2$.
Starting from $(0,0)$, one checks that
\[
(0,0), (1,0), \dots, (n-1,0), (0,1), \dots, (n-1,n-1)
\]
is a single orbit of length $n^2$.
Hence $\alpha$ is a directed Hamilton cycle.

Now let $\sigma(u,v)=(v,u)$.
Then $\beta=\sigma\circ\alpha\circ\sigma^{-1}$, so $\beta$ is conjugate to
$\alpha$ and is also a directed Hamilton cycle.
Therefore $\Cdir_n\square\Cdir_n$ decomposes into two directed Hamilton cycles.
\end{proof}

\begin{corollary}[Self-square lift]\label{cor:self-square}
Let $G$ be a finite digraph admitting a Hamilton decomposition.
Then $G\square G$ admits a Hamilton decomposition as well.
More precisely, if $G$ decomposes into $k$ directed Hamilton cycles, then
$G\square G$ decomposes into $2k$ directed Hamilton cycles.
\end{corollary}

\begin{proof}
Apply Theorem~\ref{thm:power-lift} with $b=2$, using
Lemma~\ref{lem:square}.
\end{proof}

\section{Directed tori and composite expansion}

For integers $d\ge 1$ and $m\ge 2$, define
\[
\torus{d}{m} := \Cay\bigl((\mathbb Z/m\mathbb Z)^d,\{e_0,e_1,\dots,e_{d-1}\}\bigr).
\]
Equivalently,
\[
\torus{d}{m}\cong \underbrace{\Cdir_m\square\cdots\square\Cdir_m}_{d\text{ factors}}.
\]

\begin{theorem}[Pointwise composite expansion]\label{thm:pointwise-composite}
Let $a,b\ge 2$ and $m\ge 2$.
Assume that $\torus{a}{m}$ admits a Hamilton decomposition and that
$\torus{b}{m^a}$ admits a Hamilton decomposition.
Then $\torus{ab}{m}$ admits a Hamilton decomposition.
\end{theorem}

\begin{proof}
Set
\[
G:=\torus{a}{m}.
\]
Then $G$ has
\[
|V(G)|=m^a=:n
\]
vertices, and by hypothesis its arc set is partitioned by $a$ directed
Hamilton cycles.
Also,
\[
\torus{b}{n}=\torus{b}{m^a}\cong \Cdir_n^{\square b}
\]
admits a Hamilton decomposition by assumption.
Therefore Theorem~\ref{thm:power-lift} applies to $G$ and gives a Hamilton
decomposition of $G^{\square b}$.

Finally,
\[
G^{\square b}
\cong
\bigl(\Cdir_m^{\square a}\bigr)^{\square b}
\cong
\Cdir_m^{\square ab}
\cong
\torus{ab}{m}.
\]
Hence $\torus{ab}{m}$ admits a Hamilton decomposition.
\end{proof}

\begin{corollary}[Multiplicative closure across all moduli]
\label{cor:multiplicative-closure}
Let $a,b\ge 2$.
Assume that for every integer $n\ge 3$, both $\torus{a}{n}$ and $\torus{b}{n}$
admit Hamilton decompositions.
Then for every integer $m\ge 3$, the directed torus $\torus{ab}{m}$ admits a
Hamilton decomposition.
\end{corollary}

\begin{proof}
Fix $m\ge 3$.
By hypothesis, $\torus{a}{m}$ admits a Hamilton decomposition.
Also $m^a\ge 3$, so the same hypothesis gives a Hamilton decomposition of
$\torus{b}{m^a}$.
Now apply Theorem~\ref{thm:pointwise-composite}.
\end{proof}

\begin{corollary}[Prime reduction]
\label{cor:prime-reduction}
Assume that for every prime $p$ and every integer $n\ge 3$, the directed torus
$\torus{p}{n}$ admits a Hamilton decomposition.
Then for every integer $d\ge 2$ and every integer $m\ge 3$, the directed torus
$\torus{d}{m}$ admits a Hamilton decomposition.
\end{corollary}

\begin{proof}
Proceed by induction on $d$.
If $d$ is prime, there is nothing to prove.
If $d$ is composite, write $d=ab$ with $2\le a,b<d$.
By the induction hypothesis, both families $\torus{a}{n}$ and $\torus{b}{n}$
admit Hamilton decompositions for every $n\ge 3$.
Corollary~\ref{cor:multiplicative-closure} then gives the result for $d=ab$.
\end{proof}

\begin{corollary}[Immediate consequences of currently solved cases]
\label{cor:examples}
The following implications are formal consequences of the product theorem.
\begin{enumerate}
    \item Since $\torus{2}{n}$ is solved by Lemma~\ref{lem:square}, any solved
    dimension $d$ automatically gives the doubled dimension $2d$.
    \item If the known $d=3$ theorem is available for all $n\ge 3$, then the
    dimensions $6$, $9$, $12$, $18$, and more generally every
    $2^{\alpha}3^{\beta}$ with $\alpha+\beta\ge 1$, follow by repeated
    composite expansion.
    \item If the known $d=4$ theorem is available for all $n\ge 3$, then the
    dimensions $8$, $12$, $16$, $24$, and every other dimension obtained by
    multiplying together already-solved dimensions, follow as well.
\end{enumerate}
\end{corollary}

\begin{remark}
This is a theorem-level reduction rather than a local witness construction.
It proves existence in composite dimensions once the relevant lower-dimensional
inputs are known, but it does not itself yield the kind of small low-layer rule
table used in explicit constructions for dimensions such as $3$ and $4$.
\end{remark}

\end{document}
